\section{关系}

\subsection{$n$ 元关系}

\begin{definition}[关系]
    设 $A_1, A_2, \dots, A_n$ 为 $n$ 个任意的集合，$R \subseteq A_1 \times A_2 \times \dots \times A_n$，那么
    \begin{itemize}
        \item 称 $R$ 为 $A_1, A_2, \dots, A_n$ 间的 $n$ 元\textbf{关系}；
        \item 若 $n = 2$，称 $R$ 为 \textbf{$A_1$ 到 $A_2$ 的二元关系}；
        \item 若 $R = \varnothing$，称 $R$ 为空关系；
        \item 若 $R = A_1 \times A_2 \times \dots \times A_n$，称 $R$ 为\textbf{全关系}；
        \item 若 $A_1 = A_2 = \dots = A_n = A$，称 $R$ 为 \textbf{$A$ 上的 $n$ 元关系}。
    \end{itemize}
\end{definition}

对于二元关系 $R$，可以记做：
\begin{itemize}
    \item $\langle x, y \rangle \in R \iff x \, R \, y$；
    \item $\langle x, y \rangle \notin R \iff x \, \overline{R} \, y$。
\end{itemize}

\begin{definition}[关系的相等]
    设 $R_1$ 是 $A_1, A_2, \dots, A_n$ 上的 $n$ 元关系，$R_2$ 是 $B_1, B_2, \dots, B_m$ 上的 $m$ 元关系，那么 $R_1$ 和 $R_2$ 相等当且仅当：
    \begin{itemize}
        \item $n = m$；
        \item $A_i = B_i\ (1 \le i \le n)$；
        \item 从集合来看，$R_1 = R_2$。
    \end{itemize}
\end{definition}

\subsection{二元关系的表示}

\subsubsection{图表示}

对于一个 $X$ 到 $Y$ 的二元关系 $R$，建立一张图 $G = (V, E)$：

\begin{itemize}
    \item $V = X \cup Y$；
    \item $E = \{(x, y): x \, R \, y\}$。
\end{itemize}

这个图称为关系 $R$ 的\textbf{图表示}。

\subsubsection{矩阵表示}

二元关系 $R$ 的图表示的邻接矩阵称为 $R$ 的\textbf{矩阵表示}。

\subsection{二元关系的性质}

\begin{definition}[自反性]
    设 $R$ 是集合 $X$ 上的二元关系，那么称 $R$ 是\textbf{自反的}当且仅当
    $$
    \forall\,x \in X: \langle x, x \rangle \notin R
    $$
\end{definition}

在图表示中，自反的关系没有自环。

\begin{definition}[对称性]
    设 $R$ 是集合 $X$ 上的二元关系，那么称 $R$ 是\textbf{对称的}当且仅当
    $$
    \forall\,x, y \in X: \langle x, y \rangle \in R \iff \langle y, x \rangle \in R
    $$
\end{definition}

满足对称性的关系，在图表示中边都是双向的；在矩阵表示中矩阵是对称矩阵。

\begin{definition}[传递性]
    设 $R$ 是集合 $X$ 上的二元关系，那么称 $R$ 是\textbf{传递的}当且仅当
    $$
    \forall\,x, y, z \in X: \langle x, y \rangle, \langle y, z \rangle \in R \implies \langle x, z \rangle \in R
    $$
\end{definition}

满足传递性的关系，其图表示是一个传递闭包。
